On place dans un ballon de baudruche $\qty{5.3}{\g}$ de carbonate de sodium
solide \ch{Na2CO3} auquel on ajoute un volume $V=\qty{50}{\mL}$ d'une solution
d'acide chlorhydrique de concentration $c=\qty{0.8}{\mol\per\L}$. Il se forme du
dioxyde de carbone, de l'eau et des ions sodium en solution.
\begin{enumerate}
    \item Écrire l'équation de la réaction.
    \item Quel est le réactif limitant~?
    \item Que vaut l'avancement maximal~?
    \item Décrire le système dans l'état final.
    \item Quel est le volume de gaz obtenu~? (On suppose la pression égale à
          l'intérieur et à l'extérieur du ballon.)
\end{enumerate}
\begin{donnees}
    Volume molaire~: $V_{m}=\qty{24}{\L\per\mol}$.
\end{donnees}

